\section*{Chapitre \ref{ch:rc:bermudans-and-callables} --- Bermudéennes et remboursables par anticipation}

\begin{solution}{exo:rc:bermudans-and-callables:1}
Borne inférieure~: l'européenne à trois ans, 249.5 points de base.
Borne supérieure~: la somme des neuf, 1\,647.5. La \omterm{def:rc:bermudans-and-callables:bermudan}{swaption bermudéenne}, 358.4, se situe entre les deux, bien plus près de la borne inférieure.
\end{solution}

\begin{solution}{exo:rc:bermudans-and-callables:2}
La valeur d'une européenne vaut à peu près l'annuité, multipliée par la volatilité et par la racine carrée de l'échéance~: l'option
à un an porte le swap le plus long (annuité la plus grande), mais laisse peu de temps
aux taux pour bouger~; celle à neuf ans dispose du temps, mais d'un swap d'un an. Le produit culmine
ici à trois ans.
\end{solution}

\begin{solution}{exo:rc:bermudans-and-callables:3}
Elle a vendu à l'émetteur une option bermudéenne de remboursement à la valeur capitalisée, une
option de type receveur sur les taux longs (l'émetteur rembourse lorsque les taux baissent). L'émetteur
la transmet, via le swap, au dealer, qui la couvre en vendant de la volatilité de swaptions sur le
marché.
\end{solution}

\begin{solution}{exo:rc:bermudans-and-callables:4}
Avec moins de dates d'exercice restantes, attendre vaut moins, si bien que l'exercice exige un
mouvement plus faible en faveur du détenteur~; à la dernière date, la bermudéenne se réduit à une européenne et
l'exercice a lieu dès que le swap est dans la monnaie. La frontière d'exercice reste sous le forward, car à la
monnaie la valeur de continuation (option sur les dates ultérieures)
dépasse la valeur intrinsèque immédiate.
\end{solution}

\begin{solution}{exo:rc:bermudans-and-callables:5}
La meilleure européenne reste fixe, car on recalibre $\sigma$ sur la swaption cinq ans dans cinq ans, si bien qu'elle bouge
peu. Un retour à la moyenne plus fort rend moins corrélés les taux de swap de dates d'exercice différentes~: un taux bas
à l'année trois en dit moins sur l'année six, si bien que le droit de choisir la date plus tard vaut davantage
(101.9, 109.0 et 120.0 points de base).
\end{solution}

\begin{solution}{exo:rc:bermudans-and-callables:6}
Ajustée et appliquée sur les mêmes trajectoires, la régression utilise le futur propre de chaque trajectoire pour
décider s'il faut exercer, un biais d'anticipation qui tire le prix vers le haut~; appliquée à des trajectoires
indépendantes, la règle est une règle d'arrêt légitime (sous-optimale) dont le gain espéré constitue une borne
inférieure.
\end{solution}

\begin{solution}{exo:rc:bermudans-and-callables:7}
L'obligation zéro-coupon remboursable par anticipation (callable) vaut le pair à un rendement de capitalisation de 5.61~\%~; le
zéro-coupon ordinaire de trente ans sur la courbe rapporte 4.13~\% (annuel). L'investisseur gagne 1.48
point de pourcentage par an pour la clause de remboursement anticipé.
\end{solution}

\begin{solution}{exo:rc:bermudans-and-callables:8}
L'européenne ne couvre que la borne inférieure~: l'\omterm{def:rc:bermudans-and-callables:switch}{option d'échange} (ici 44~\% de la
meilleure européenne) n'est pas couverte, et le vega de la bermudéenne se répartit sur l'ensemble des \omterm{def:rc:bermudans-and-callables:coterminal}{swaptions co-terminales},
se déplaçant avec les taux~; le notionnel d'une couverture neutre en delta et en vega n'est pas celui de la bermudéenne.
Il faut couvrir le vega et le delta de chaque \omterm{def:rc:bermudans-and-callables:coterminal}{swaption co-terminale}, et
constituer une réserve pour provisionner l'\omterm{def:rc:bermudans-and-callables:switch}{option d'échange}.
\end{solution}

\begin{solution}{pb:rc:bermudans-and-callables:1}
\textbf{1.} 4.13~\%.
\textbf{2.} 5.61~\%.
\textbf{3.} 1.526.
\textbf{4.} 0.526 par unité émise~: 1.526 moins le pair.
\textbf{5.} Pour 1.48 point de pourcentage par an de rendement supplémentaire, et parce que la
réglementation a fait de ces obligations le moyen le plus efficace de détenir une duration longue en dollars.
\textbf{6.} Cinquième année, première date de remboursement anticipé~: 0.456.
\textbf{7.} 0.070 par unité.
\textbf{8.} Une obligation zéro-coupon remboursée à sa valeur capitalisée est dans la monnaie au plus haut dès le début
(la capitalisation croît au-dessus des taux de la courbe dès l'origine), et le premier remboursement anticipé
réduit le passif restant le plus long.
\textbf{9.} Une option bermudéenne de type receveur sur les taux longs, reçue via le swap de l'émetteur, contre
laquelle le dealer verse le financement à taux variable.
\textbf{10.} C'est une large position longue en volatilité sur des options peu liquides de longue échéance~;
les dealers sont rémunérés pour porter et couvrir cette position, non pour la conserver.
\textbf{11.} 2.33~millions de dollars par point de base de $\sigma$ du modèle et par milliard de dollars de notionnel.
\textbf{12.} Des swaptions de longue échéance (co-terminales de swaps de trente ans à partir de cinq ans), et
des swaptions receveuses d'échéance plus longue.
\textbf{13.} Environ 23.3~millions de dollars par point de base de $\sigma$.
\textbf{14.} Elle est tirée vers le bas par les ventes régulières~; les dealers
longs des bermudéennes les valorisent plus bas à mesure que la volatilité baisse.
\textbf{15.} La bermudéenne disparaît lorsqu'elle est remboursée~: le dealer se retrouve à découvert
de la volatilité qu'il a vendue et doit la racheter, souvent dans un marché de taux en baisse et de volatilité
en hausse.
\textbf{16.} Risque de réinvestissement en cas de remboursement anticipé (dans les états de taux bas), risque de change face à son passif, et
concentration sur peu d'émetteurs.
\textbf{17.} Peu de dealers peuvent porter et couvrir ces options~; leurs couvertures sont corrélées, si bien que les flux
font bouger le marché, et un renversement (les remboursements anticipés) les frappe tous
à la fois.
\textbf{18.} La structure de volatilité et de corrélation du modèle (la valeur de l'\omterm{def:rc:bermudans-and-callables:switch}{option d'échange}), le smile
des swaptions de longue échéance, et la règle d'exercice.
\textbf{19.} Résultat à retenir~: \emph{couverture Formosa}~: l'obligation zéro-coupon remboursable par anticipation est au pair
pour une capitalisation de 5.61~\% contre 4.13~\% pour le zéro-coupon ordinaire~; l'option bermudéenne de remboursement anticipé
vaut 0.070 de plus que la meilleure européenne (0.456)~; son vega est de 2.33~millions de dollars par point de base
et par milliard de dollars de notionnel.
\textbf{20.} Le dealer vend de la volatilité, parce qu'il a acheté une bermudéenne qu'il ne veut pas
porter, à un émetteur qui l'avait prise à un investisseur rémunéré pour s'en défaire.
\end{solution}

\begin{solution}{iq:rc:bermudans-and-callables:1}
Le droit d'entrer, à l'une quelconque de plusieurs dates, dans un swap qui se termine à une date fixe. Il
vaut au moins autant qu'une européenne seule (engagée à cette date), et davantage, car le
détenteur choisit la date après avoir observé les taux~: l'\omterm{def:rc:bermudans-and-callables:switch}{option d'échange}.
\iqlookfor{la définition et l'\omterm{def:rc:bermudans-and-callables:switch}{option d'échange}.}
\end{solution}

\begin{solution}{iq:rc:bermudans-and-callables:2}
Simuler l'état~; en remontant depuis la dernière date d'exercice, régresser les flux futurs réalisés actualisés des
trajectoires dans la monnaie sur des fonctions de base de l'état.
et exercer là où la valeur d'exercice dépasse la continuation ajustée. Toute
règle d'exercice fournit un gain réalisable, si bien que sa valeur (sur des trajectoires indépendantes de l'ajustement)
est au plus égale à l'optimum.
\iqlookfor{la régression rétrograde et la raison pour laquelle elle fournit une borne inférieure.}
\end{solution}

\begin{solution}{iq:rc:bermudans-and-callables:3}
Retour à la moyenne~: il décorrèle les taux de swap des dates d'exercice successives~; avec
les prix des swaptions européennes étant maintenus en recalibrant $\sigma$, un $\kappa$ plus élevé donne une
\omterm{def:rc:bermudans-and-callables:switch}{option d'échange} plus grande.
\iqlookfor{le retour à la moyenne vu comme corrélation entre les dates d'exercice.}
\end{solution}

\begin{solution}{iq:rc:bermudans-and-callables:4}
Les émetteurs d'obligations remboursables par anticipation les échangent contre des swaps auprès de dealers et répercutent le call~; les dealers,
détenant des \omterm{def:rc:bermudans-and-callables:bermudan}{swaptions bermudéennes} longues, se couvrent en vendant des swaptions à long terme, une offre régulière qui
pèse sur la volatilité implicite à long terme~; les remboursements par anticipation ou les dénouements inversent ce flux.
\iqlookfor{la chaîne allant de l'investisseur au marché des swaptions.}
\end{solution}

\begin{solution}{iq:rc:bermudans-and-callables:5}
Comparer aux swaptions européennes (limite à une seule date) et aux bornes~; arbre contre régression sur
le même modèle~; convergence en pas et en trajectoires~; une borne supérieure duale pour l'écart~;
sensibilité aux paramètres du modèle ($\kappa$, $\sigma(t)$, corrélation) et au
ensemble de calibration~; stabilité des grecques~; comparaison à un modèle indépendant pris comme référence.
\iqlookfor{un plan de validation structuré.}
\end{solution}

\begin{solution}{iq:rc:bermudans-and-callables:6}
Avec une erreur type d'environ deux points de base, un mouvement de deux points de base est
du bruit, et non un bogue. Indiquer l'erreur type~; utiliser davantage de trajectoires, des variables antithétiques ou
des variables de contrôle (les européennes ont des formes fermées), fixer la graine et le
jeu de trajectoires pour les grecques (nombres aléatoires communs), et comparer à l'arbre.
\iqlookfor{erreurs types et réduction de variance.}
\end{solution}
